\section{Actual smoothness of the original variety and parameter family}

The original chart condition is algebraic smoothness of the original
quotient chart algebras. This section proves smoothness of the actual
native scheme morphisms. The whole defining ideal, original field action,
original chart coordinates, and original parameter projection are retained.
No scheme comparison or smoothness certificate is an extra input.

\begin{lemma}[The original affine coordinate rings coincide as algebras]
\label{lem:c0n-family-original-chart-coordinates}
\lean{EnvelopeProperty.Char0.MultiplierFamilySmoothness.chartCoordinateEquiv,
EnvelopeProperty.Char0.MultiplierFamilySmoothness.rename_dehom,
EnvelopeProperty.Char0.MultiplierFamilySmoothness.chartCoordinateEquiv_chartMk}
\leanok
For every original coordinate $j$, the old degree-zero chart algebra
$k[x]_{(x_j)}$ is isomorphic as a $k$-algebra to the native affine coordinate
polynomial ring with variables indexed by $i\ne j$. On a homogeneous
original polynomial $F$ of degree $d$, the actual class $F/x_j^d$ goes to
its original dehomogenization at $x_j=1$.
\end{lemma}
\begin{proof}
Use the proved original chart dehomogenization isomorphism to the
polynomial ring in $\dim\mathbb P$ variables, followed by the actual
renaming equivalence $a\mapsto j.\mathrm{succAbove}(a)$ onto the variables
other than $j$. Checking the selected coordinate and each remaining
coordinate identifies the entire polynomial homomorphism with the native
dehomogenization. The original homogeneous-fraction formula then gives
the claimed identity on every original form, including its coefficient
action.
\end{proof}

\begin{lemma}[Whole native chart sections are the original quotient chart algebra]
\label{lem:c0n-family-original-chart-quotient}
\lean{EnvelopeProperty.Char0.MultiplierFamilySmoothness.chartIdeal_map_eq,
EnvelopeProperty.Char0.MultiplierFamilySmoothness.coordinateQuotientEquiv,
EnvelopeProperty.Char0.MultiplierFamilySmoothness.nativeChartOpen,
EnvelopeProperty.Char0.MultiplierFamilySmoothness.nativeChartOpen_isAffine,
EnvelopeProperty.Char0.MultiplierFamilySmoothness.nativeChartAlgebra,
EnvelopeProperty.Char0.MultiplierFamilySmoothness.nativeChartCoordinateSectionsEquiv,
EnvelopeProperty.Char0.MultiplierFamilySmoothness.nativeChartSectionsEquiv,
EnvelopeProperty.Char0.MultiplierFamilySmoothness.nativeChartCoordinateSectionsEquiv_pullback,
EnvelopeProperty.Char0.MultiplierFamilySmoothness.coordinateQuotientEquiv_mk_homogeneous,
EnvelopeProperty.Char0.MultiplierFamilySmoothness.nativeChartSectionsEquiv_pullback_homogeneous,
EnvelopeProperty.Char0.MultiplierFamilySmoothness.nativeChartSectionsEquiv_coefficients}
\leanok
\uses{lem:c0n-family-original-chart-coordinates}
Let $j$ be an original coordinate not belonging to the original prime
homogeneous ideal $I$. The actual native inverse-image coordinate open
has its entire section algebra isomorphic to the original quotient chart
algebra $k[x]_{(x_j)}/I_{(x_j)}$, with every original dehomogenized
polynomial and scalar preserved. This is an equality of whole ideal
quotients, not an equality of supports or radicals.
\end{lemma}
\begin{proof}
For one inclusion, write any element of the original chart ideal as a
homogeneous fraction $F/x_j^d$ with $F\in I$, using the actual prime ideal
and $x_j\notin I$. The preceding coordinate equivalence sends it to the
original dehomogenization of $F$, so it belongs to the entire native
image ideal. Conversely decompose any original $F\in I$ into its finitely
many homogeneous components, all still in $I$. The original chart
fractions of these components belong to $I_{(x_j)}$ and map to their
actual dehomogenizations; adding them gives the dehomogenization of $F$.
Thus the two entire ideals correspond. Apply the genuine quotient
algebra equivalence and compose with the proved native quotient-section
algebra equivalence. All maps preserve the original scalar action by
construction.
\end{proof}

\begin{theorem}[Original smooth algebra charts imply smoothness of the original native field map]
\label{thm:c0n-original-native-variety-smooth}
\lean{EnvelopeProperty.Char0.MultiplierFamilySmoothness.nativeChartSections_smooth,
EnvelopeProperty.Char0.MultiplierFamilySmoothness.nativeChart_structure,
EnvelopeProperty.Char0.MultiplierFamilySmoothness.nativeChart_smooth,
EnvelopeProperty.Char0.MultiplierFamilySmoothness.nativeAffineCover,
EnvelopeProperty.Char0.MultiplierFamilySmoothness.nativeStructure_smooth}
\leanok
\uses{lem:c0n-family-original-chart-quotient}
If the original projective subvariety has smooth charts, its actual
native structure morphism to the original field spectrum is smooth.
No native scheme smoothness, chosen atlas or field-map comparison is
supplied.
\end{theorem}
\begin{proof}
The inverse images of the actual ambient coordinate opens cover the
same original native closed scheme. On a coordinate not in $I$, transport
the given original quotient-algebra smoothness along the preceding whole
algebra equivalence. The native affine-spectrum smoothness criterion and
the actual affine-open spectrum isomorphism prove smoothness of that
same original chart field morphism. Its scalar map is the original
structure-map action from the section-algebra comparison. If $x_j\in I$,
its native dehomogenization is $1$, so the entire coordinate ideal is the
unit ideal. Its actual coordinate ring is the zero ring, which is
standard smooth of relative dimension zero over the field; thus this
empty chart is included correctly. Native source-locality for smoothness
on the actual coordinate cover proves the displayed original global
structure morphism smooth.
\end{proof}

\begin{lemma}[The whole actual parameter quotient is the original variety product]
\label{lem:c0n-family-original-product-iso}
\lean{EnvelopeProperty.Char0.MultiplierFamilySmoothness.familyProductInclusion,
EnvelopeProperty.Char0.MultiplierFamilySmoothness.familyProductInclusion_isClosedImmersion,
EnvelopeProperty.Char0.MultiplierFamilySmoothness.familyProduct_kernel,
EnvelopeProperty.Char0.MultiplierFamilySmoothness.familyProductIso,
EnvelopeProperty.Char0.MultiplierFamilySmoothness.familyProductIso_inclusion,
EnvelopeProperty.Char0.MultiplierFamilySmoothness.familyProductIso_projection,
EnvelopeProperty.Char0.MultiplierFamilySmoothness.familyPullbackIso,
EnvelopeProperty.Char0.MultiplierFamilySmoothness.familyPullbackIso_projection,
EnvelopeProperty.Char0.MultiplierFamilySmoothness.familyToOriginal}
\leanok
The same original projective quotient obtained by extending $I$ to
$k[t][x]$ is isomorphic to the actual original closed product model of
$X\times\mathbb A^1$. Its entire original polynomial-base projection
commutes with the isomorphism.
\end{lemma}
\begin{proof}
Compose the actual homogeneous quotient closed immersion with the proved
ambient polynomial-line product isomorphism. Its actual kernel ideal
sheaf is the native polynomial-line ideal of the entire coefficient
extension of $I$. The proved whole original product-cut-ideal identity
identifies precisely that kernel with the original closed product model's
ideal. An actual closed immersion is isomorphic to its genuine scheme
image; transport this genuine image scheme along that whole kernel-ideal
equality. The native image-inclusion triangle and the actual ambient
product projection identity give the whole asserted original projection
triangle. No scheme isomorphism is assumed.
\end{proof}

\begin{theorem}[The same original parameter family is smooth over the parameter and over the original field]
\label{thm:c0n-actual-original-family-smooth}
\lean{EnvelopeProperty.Char0.MultiplierFamilySmoothness.nativeProductStructure_smooth,
EnvelopeProperty.Char0.MultiplierFamilySmoothness.familyProjection_smooth,
EnvelopeProperty.Char0.MultiplierFamilySmoothness.polynomialBase_smooth,
EnvelopeProperty.Char0.MultiplierFamilySmoothness.familyAbsoluteStructure,
EnvelopeProperty.Char0.MultiplierFamilySmoothness.familyAbsoluteStructure_smooth,
EnvelopeProperty.Char0.MultiplierFamilySmoothness.familyToOriginal_structure}
\leanok
\uses{thm:c0n-original-native-variety-smooth,lem:c0n-family-original-product-iso}
Under the original smooth-chart condition, the actual original parameter
projection is smooth, and its actual composite with $\operatorname{Spec}k[t]
\to\operatorname{Spec}k$ is smooth. The original whole parameter-family
scheme and both original coefficient maps are retained. Pure relative
dimension is a separate subsequent consequence of this smoothness and
the already proved integrality.
\end{theorem}
\begin{proof}
The original native closed product model is genuinely isomorphic over
$k[t]$ to the pullback of the original native variety field map along the
actual polynomial coefficient morphism. The preceding original variety
smoothness is preserved by this genuine base change. Its original
polynomial-base projection is therefore smooth. Transport through the
preceding same-family isomorphism and whole projection triangle to prove
smoothness of the same original parameter projection. The actual
polynomial algebra $k[t]$ is smooth over $k$, so composition gives the
asserted same-family total field smoothness. No dimension or vanishing
certificate is inserted in this argument.
\end{proof}
